% =====================================================================
%  Preámbulo común para los apuntes de Física Electrónica
%  Uso (desde unidades/XX-nombre/archivo.tex):
%     \documentclass[11pt,a4paper]{article}
%     \input{../../comun/preambulo.tex}
% =====================================================================

% ---------- Idioma y tipografía ----------
\usepackage[T1]{fontenc}
\usepackage[utf8]{inputenc}
\usepackage[spanish,es-nodecimaldot,es-noshorthands,es-tabla]{babel}
\usepackage{lmodern}
\usepackage{microtype}
\usepackage[a4paper,margin=2.3cm,headheight=14pt]{geometry}
\usepackage{parskip}
\setlength{\parskip}{0.55em}

% ---------- Matemática y unidades ----------
\usepackage{amsmath,amssymb,mathtools}
\usepackage{bm}
\usepackage[per-mode=symbol,output-decimal-marker={.},separate-uncertainty]{siunitx}
\sisetup{range-phrase={\text{--}},range-units=single}
\DeclareSIUnit{\eV}{eV}

% ---------- Gráficos ----------
\usepackage{graphicx}
\usepackage{xcolor}
\usepackage{tikz}
\usetikzlibrary{arrows.meta,decorations.pathmorphing,decorations.markings,
  patterns,calc,positioning,shapes.geometric,fit,backgrounds,shadings}
\usepackage{pgfplots}
\pgfplotsset{compat=1.18}
\usepgfplotslibrary{fillbetween}
\usepackage[siunitx,RPvoltages]{circuitikz}

% ---------- Tablas, listas, figuras ----------
\usepackage{booktabs,array,tabularx,multirow}
\usepackage{enumitem}
\setlist{itemsep=0.2em,topsep=0.3em}
\usepackage{caption,subcaption}
\captionsetup{font=small,labelfont=bf,format=hang}
\usepackage{float}
\usepackage{wrapfig}
\usepackage{needspace}

% ---------- Colores ----------
\definecolor{azul}{HTML}{1F4E79}
\definecolor{azulclaro}{HTML}{DCE9F5}
\definecolor{verde}{HTML}{2E7D32}
\definecolor{verdeclaro}{HTML}{E3F2E1}
\definecolor{naranja}{HTML}{C55A11}
\definecolor{naranjaclaro}{HTML}{FBE9DB}
\definecolor{rojo}{HTML}{B71C1C}
\definecolor{rojoclaro}{HTML}{FDECEA}
\definecolor{violeta}{HTML}{5E35B1}
\definecolor{violetaclaro}{HTML}{EDE7F6}
\definecolor{gris}{HTML}{555555}
% colores para dibujos de bandas / portadores
\colorlet{colEc}{azul}
\colorlet{colEv}{rojo}
\colorlet{colElec}{azul}
\colorlet{colHueco}{rojo}
\colorlet{colFoton}{naranja}

% ---------- Cajas didácticas ----------
\usepackage[most]{tcolorbox}
\tcbset{
  base/.style={enhanced,breakable,arc=2pt,boxrule=0.6pt,left=6pt,right=6pt,
    top=4pt,bottom=4pt,fonttitle=\bfseries\small,before skip=8pt,after skip=8pt}
}
% Definición
\newtcolorbox{definicion}[1][]{base,colback=azulclaro,colframe=azul,
  title={Definición\ifx\relax#1\relax\else: #1\fi}}
% Resultado / fórmula clave
\newtcolorbox{clave}[1][]{base,colback=verdeclaro,colframe=verde,
  title={\ifx\relax#1\relax Resultado clave\else #1\fi}}
% Ejemplo resuelto
\newtcolorbox{ejemplo}[1][]{base,colback=white,colframe=violeta,
  coltitle=white,colbacktitle=violeta,
  title={Ejemplo\ifx\relax#1\relax\else: #1\fi}}
% Corrección / aclaración respecto de las diapositivas
\newtcolorbox{ojo}[1][]{base,colback=rojoclaro,colframe=rojo,
  title={\ensuremath{\triangle}\ Ojo\ifx\relax#1\relax\else: #1\fi}}
% Intuición física / comentario
\newtcolorbox{intuicion}[1][]{base,colback=naranjaclaro,colframe=naranja,
  title={Idea física\ifx\relax#1\relax\else: #1\fi}}

% ---------- Encabezados ----------
\usepackage{fancyhdr}
\pagestyle{fancy}
\fancyhf{}
\renewcommand{\headrulewidth}{0.4pt}
\fancyhead[L]{\small\color{gris}Física Electrónica}
\fancyhead[R]{\small\color{gris}\unidadtitulo}
\fancyfoot[C]{\small\thepage}
\newcommand{\unidadtitulo}{}
\newcommand{\unidad}[1]{\renewcommand{\unidadtitulo}{#1}}

\usepackage{titlesec}
\titleformat{\section}{\Large\bfseries\color{azul}}{\thesection}{0.6em}{}
\titleformat{\subsection}{\large\bfseries\color{azul}}{\thesubsection}{0.6em}{}
\titleformat{\subsubsection}{\normalsize\bfseries\color{azul!80!black}}{\thesubsubsection}{0.6em}{}

\usepackage[hidelinks,bookmarksnumbered]{hyperref}
\usepackage{bookmark}

% ---------- Macros físicas ----------
\newcommand{\EG}{E_G}
\newcommand{\Ec}{E_c}
\newcommand{\Ev}{E_v}
\newcommand{\EF}{E_F}
\newcommand{\Eph}{E_{\mathrm{ph}}}
\newcommand{\lG}{\lambda_G}
\newcommand{\VT}{V_T}
\newcommand{\Popt}{P_{\mathrm{opt}}}
\newcommand{\Resp}{\mathcal{R}}
\newcommand{\vsat}{v_{\mathrm{sat}}}
\newcommand{\Voc}{V_{oc}}
\newcommand{\Isc}{I_{sc}}
\newcommand{\diapo}[1]{{\small\textsf{[diapo~#1]}}}
\newcommand{\dd}{\,\mathrm{d}}

% Distintivo para resultados que ya están en la hoja de fórmulas oficial.
% Uso: en el título de una caja, \begin{clave}[Título \enhoja], o en texto con \enhojatexto.
\newcommand{\enhojatexto}{{\scriptsize\sffamily\colorbox{verdeclaro}{\textcolor{verde!70!black}{\ensuremath{\checkmark}\,en la hoja}}}}
\newcommand{\enhoja}{\hfill{\scriptsize\sffamily\colorbox{white}{\textcolor{verde!70!black}{\ensuremath{\checkmark}\,en la hoja}}}}

\input{../../comun/constantes.tex}

% ---------- Estilos TikZ comunes ----------
\tikzset{
  banda/.style={line width=1.1pt},
  bandac/.style={banda,colEc},
  bandav/.style={banda,colEv},
  fermi/.style={dashed,gris,line width=0.6pt},
  elec/.style={circle,fill=colElec,inner sep=0pt,minimum size=5pt},
  hueco/.style={circle,draw=colHueco,fill=white,line width=0.9pt,inner sep=0pt,minimum size=5pt},
  foton/.style={colFoton,line width=0.9pt,decorate,
    decoration={snake,amplitude=1.6pt,segment length=5pt,post length=3pt},-{Stealth[length=5pt]}},
  flecha/.style={-{Stealth[length=5pt]},line width=0.7pt},
}
